\providecommand{\toplevelprefix}{../..}  %
\documentclass[../../book-main_zh.tex]{subfiles}

\begin{document}

\chapter*{前言}

%\vspace{-2cm}
\begin{center}
    ``{\em 条条大路通罗马}。''

\end{center}
\vspace{5mm}

\paragraph{目的。}
本书揭示并深入研究了几乎所有现代机器（人工）智能实践背后一个普遍且根本的问题。
即{\em 如何在高维空间中有效且高效地学习数据的低维分布，
    然后将该分布转化为紧凑且结构化的表征？} 对于任何智能系统，无论是自然的还是人造的，
这种表征通常可以被视为从外部世界感知的数据中学习到的{\em 记忆}或{\em 经验知识}。近年来，这些知识也常被非正式的称为“世界模型”。

\paragraph{对象。}
本教科书旨在为{\em 高年级本科生和一年级研究生}系统地介绍学习此类数据分布的（深度）
表征的数学和计算原理。本书的主要先修知识包括本科阶段的线性代数、
概率论/统计学和优化算法。如果对信号处理（特别是稀疏表征和压缩感知）、信息论以及
控制理论的基本概念有所了解，将有助于更好地理解本书内容。

\paragraph{动机。}
撰写本书的主要动机是，在过去几年中，作者及众多同行在通过严谨的方法论来建立起能理解深度神经网络乃至智能本身的数学理论方面取得了重要进展。
这种新方法所倡导的演绎方法论，与当前人工智能实践中占主导地位的主要基于归纳和试错的方法论形成了鲜明对比，同时又高度互补。
近年来，由于人们对目前看似强大的AI模型和系统缺乏理解，导致了社会上日益增长的炒作和恐惧。
我们认为，现在比以往任何时候都更需要认真尝试建立一种科学严谨的理论框架来理解智能。
本书的总体目标是通过提供坚实的理论和实验证据，证明如今已经可以将智能作为一门科学与数学学科来研究。
我们将论证，智能的基本能力在于创造新的记忆（或知识）或修正已有记忆（或知识）。因此，正如本书副标题所示，
本书可被视为智能在学习经验性知识或者记忆的层面上，对建立一套{\em 智能的数学理论}的首次尝试。

在技术层面上，本书提出的理论框架有助于弥合长期存在的鸿沟：一方是主要基于解析的几何、
代数和概率模型（例如子空间、高斯分布和方程）来对数据结构进行建模的经典方法，
另一方是基于经验设计的非参数的、数据驱动模型（例如深度网络）的“现代”方法。
事实证明，如果人们认识到这两种看似不同的方法论都在试图对感兴趣的数据分布中的
{\em 低维}结构进行学习与表征，那么它们的统一就成为可能，甚至顺理成章。
它们其实都是在寻求、表示和利用数据分布的低维结构的不同方式。从这个角度来看，即使是许多看似无关的、
在不同时期于不同领域独立开发的计算技术，现在就可以在一个共同的理论和计算框架下得到更好的理解，
并且从现在起可能可以放在一起进行研究和完善。正如我们将在本书中看到的，这些技术包括但不限于：
\begin{enumerate}
    \item 经典信息论和编码理论中发展的有损压缩编码与解码，和现代表征学习；
    \item 经典信号处理中的去噪和降维与现代生成方法的扩散和去噪模型；
    \item 通过最大后验估计或条件采样进行的贝叶斯推断与通过连续化技术\footnote{例如增广拉格朗日方法}的约束优化。
\end{enumerate}

\paragraph{内容。}
我们相信，本书提出的统一概念和计算框架，对于那些真正希望澄清关于深度神经网络和
（人工）智能的谜团与误解的读者来说，将具有巨大的价值。此外，该框架旨在为读者提供指导原则，
以便在未来开发出更好、真正智能的系统。更具体地说，除了非正式的引言（章）之外，
本书的主要技术内容将组织为六个密切相关的主题（章）：
\begin{enumerate}
    \item 我们将从主成分分析（PCA）、独立成分分析（ICA）和字典学习（DL）等经典且最基本的模型开始，
        这些模型假设感兴趣的低维分布具有线性和独立的结构。从这些在信号处理和压缩感知中
        得到充分研究和理解的简单且理想化模型出发，我们将介绍关于如何学习和表示
        低维分布的最基本和最重要的思想\footnote{例如去噪和降维}。

    \item 为了将这些经典模型及其求解方法推广到可能无法通过任何简单的解析模型加以表示的低维分布，
        我们引入了一个用于学习此类分布的通用计算原理：{\em 压缩}。正如我们将看到的，数据压缩
        为所有看似不同的经典和现代分布或表征学习方法提供了一个统一的视角，包括熵最小化、
        基于分数匹配的去噪、基于率失真理论的有损压缩，以及通过最大化信息增益的判别式表征学习。

    \item 在这个统一的框架内，现代深度神经网络（DNN），如ResNet、CNN和Transformer，
        都可以在数学上被解释为（展开的）优化算法，这些算法通过降低编码长度/速率或提升信息增益，
        迭代地实现更优的压缩与表征。该框架不仅有助于解释迄今为止凭经验设计的深度网络架构，
        还能启发设计出更简洁且更高效的新架构。

    \item 此外，为了确保学习到的数据分布表征是正确且一致的，包含编码和解码的
        {\em 自编码}架构变得必不可少。为了使学习系统能够完全自动化和连续化，我们将引入一个
        强大的{\em 闭环转录框架}，该框架使自编码系统能够通过编码器和解码器之间的极小极大博弈
        进行自我纠正，从而实现自我改进。

    \item 为了将理论与实践相结合，我们随后将研究如何利用学习到的数据分布和表征
        作为强大的先验，来进行贝叶斯推断或约束优化。这些推断与优化过程几乎涵盖了现代人工智能实践中的各类流行任务与应用场景，包括对图像和文本等真实世界高维数据的条件估计、补全和生成。

    \item 最后同样重要的是，我们将逐步演示如何利用大规模真实世界数据集（包括视觉数据、身体
        运动数据和文本数据）有效且高效地学习低维数据分布的深度表征，并将其应用于诸多实际应用中，
        如图像分类、图像补全、图像分割、图像生成、运动估计以及文本数据的类似任务。
\end{enumerate}

\paragraph{总结。}
总而言之，本书提出的技术内容在经典解析方法与现代数据驱动方法之间、简单参数模型与深度非参数模型之间、
多样化的归纳实践与基于第一性原理的统一演绎框架之间，建立起了深刻而系统的概念与技术联系。我们将揭示，许多看似无关
甚至相互竞争的方法，尽管是在不同的历史时期、在不同的领域用不同的术语发展起来的，但它们其实都在努力实现一个共同的目标：
\begin{quote}
    \centering{\em 寻求并利用嵌入在高维空间中的数据分布的内在低维结构。}
\end{quote}
为此，本书将带领读者经历一段完整的探索历程：从理论建构出发，经过严密的数学推导，提出可以通过计算实现的算法，并最终落地于实际应用。



\end{document}